\documentclass[11pt,a4paper]{article}
\usepackage[margin=2cm]{geometry}
\usepackage[T1]{fontenc}
\usepackage{lmodern}
\usepackage[english]{babel}
\usepackage{xcolor,booktabs,array,enumitem,titlesec,fancyhdr}
\usepackage[hidelinks]{hyperref}
\definecolor{navy}{HTML}{16324F}
\definecolor{okgreen}{HTML}{1B7A4E}
\definecolor{limred}{HTML}{A32020}
\titleformat{\section}{\large\bfseries\color{navy}}{\thesection}{0.6em}{}
\titleformat{\subsection}{\normalsize\bfseries\color{navy}}{\thesubsection}{0.6em}{}
\setlength{\parindent}{0pt}\setlength{\parskip}{4pt}
\pagestyle{fancy}\fancyhf{}
\fancyhead[L]{\small\color{gray}Defence preparation -- N.~Belhadj}
\fancyhead[R]{\small\color{gray}Internal document, not part of the thesis}
\fancyfoot[C]{\small\thepage}
\newcounter{qn}
\newcommand{\Q}[1]{\refstepcounter{qn}\par\medskip\noindent{\color{navy}\bfseries Q\theqn.\ #1}\par\nopagebreak}
\newcommand{\A}{\noindent\textbf{Answer.}\ }
\newcommand{\Avoid}[1]{\par\noindent{\color{limred}\small\textbf{Avoid:} #1}}
\newcolumntype{P}[1]{>{\raggedright\arraybackslash}p{#1}}
\begin{document}

\begin{center}
{\Large\bfseries\color{navy} Defence Preparation}\\[3pt]
{\large Deep Learning-Based Prediction of Disease Severity Using Medical Imaging:\\
An Application to Diabetology}\\[4pt]
{\small Every number below is taken from the thesis (version v27); section and table
numbers refer to that version.}
\end{center}

\section{The thesis in one minute}
\begin{enumerate}[leftmargin=1.5em,itemsep=1pt]
\item \textbf{Problem.} Representative DR grading systems share three limitations: they
grade each image in isolation (L1), ignore the global topology of the vessels (L2), and
rarely model the order of the grades or measure missed referrals of severe cases (L3).
\item \textbf{Objective.} Exploit relations, topology and ordinality to improve DR grading,
tested through three pre-specified hypotheses H1, H2, H3.
\item \textbf{Framework, not a single model.} TOPORET (relations $+$ vessel topology,
transductive, H1 and H2), DGTS (inductive population graph with a feature-space
topological regulariser, external EyePACS) and DOTS (ordinal CORAL pipeline evaluated
against screening targets, external Messidor-2, H3).
\item \textbf{Results.} Evidence supporting H1 ($t_4=3.82$, $p=0.019$; corrected for fold
overlap, $p=0.015$ for the three benchmarks combined); a small statistical association for H2
($F=84.3$, $p<0.001$, $\eta^2\approx0.01$); DOTS satisfies the screening targets at the point
estimate internally (QWK $0.953$, Sev-NR $0.97\%$) and on Messidor-2 (QWK $0.931$, Sev-NR
$1.4\%$), where none of four reduced variants does.
\item \textbf{Scope.} Retrospective (Level~1) evidence, not a clinical validation; the
internal Sev-NR interval $[0.2\%,\,2.8\%]$ still crosses $2\%$. The next steps are a
patient-level and multi-seed re-evaluation, a unified model with deferral, and a
prospective multi-site reader study.
\end{enumerate}

\section{Numbers to know by heart}
\begin{center}\small
\begin{tabular}{@{}P{5.2cm}P{10.8cm}@{}}
\toprule
Item & Value \\
\midrule
DR burden & about $103$ million adults (Chapter~1) \\
H1 (TOPORET, Kaggle DR) & $+0.04$ QWK, $t_4=3.82$, $p=0.019$, $d_z=1.71$; Messidor-2 $+0.07$, APTOS $+0.04$ \\
H1 robustness & bootstrap CI $[+0.021,\,+0.059]$; sign-flip $p=1/32$; Nadeau--Bengio $t=2.45/2.63/2.25$ ($p=0.070/0.058/0.088$); Fisher $p=0.015$ \\
H2 (Kaggle DR, $35{,}126$ images) & $F=84.3$, $p<0.001$, $\eta^2\approx0.0095$; $F\approx42$ if $n$ halved; $\alpha=0.005<0.05/6$ \\
TOPORET ablation (Kaggle DR) & TDA features $+0.02$; graph with TDA in edges $+0.03$; full $+0.04$ QWK \\
TOPORET settings & $\alpha=0.6$, $\tau=0.65$, five-fold, seed $42$, image-level, transductive \\
Fused corpus & $5{,}859$ images (APTOS~2019 $3{,}662$ $+$ KaggleDR-F1 $2{,}197$); $309$ Severe, $472$ PDR \\
DGTS & CV QWK $0.941\pm0.008$; EyePACS zero-shot $0.925\pm0.008$ ($35{,}126$ images); Sev-NR on EyePACS $\approx2.1\%$; $18.3$~ms/image \\
DGTS ablation (APTOS) & graph $+0.018$ (A2$\to$A3); PH regulariser $+0.030$ (A2$\to$A4); Wilcoxon $p=0.031$ \\
DOTS internal & QWK $0.9527\pm0.0063$; Sev-NR $0.97\%$ ($3/309$) $[0.2,\,2.8]\%$; PDR-NR $1.27\%$ ($6/472$) $[0.5,\,2.7]\%$; McNemar $p=0.013$ \\
DOTS Messidor-2 & $1{,}744$ images; QWK $0.931$; Sev-NR $1.4\%$; PDR-NR $1.9\%$; no $0\leftrightarrow4$ confusion \\
H3b variants (Messidor-2) & Sev-NR: baseline $5.8\%$, $+$GraphSAGE $4.1\%$, $+$CORAL $3.7\%$, $+$TTA $2.9\%$ \\
DOTS errors & $82.9\%$ adjacent; none spans four grades; about $75$~ms/image (TTA $\times5$) \\
Sample size for the next study & about $650$ severe images for the CI upper bound $<2\%$; about $1{,}240$ for $80\%$ power \\
Publications & $13$: $9$ accepted or published ($7$ conference, $2$ book chapters) $+$ $4$ journals ($1$ under revision at PLOS ONE, $3$ under review) \\
\bottomrule
\end{tabular}
\end{center}

\section{Likely questions and answers}

\subsection{Coherence and originality}

\Q{What is the single contribution of this thesis?}
\A A structured body of retrospective evidence on how relational, topological and ordinal
information can complement deep visual representations for DR grading, obtained with
three dedicated systems and pre-specified tests, and reported with its limits.

\Q{Why three systems instead of one model?}
\A The thesis proposes a three-axis research framework. Each axis is instantiated by a
dedicated system so that its contribution can be measured separately, with its own
ablation. Combining the three in one model -- a population graph with vascular topology
and a CORAL head, plus deferral -- is the first model to build next, not a claim of this thesis.
\Avoid{``DOTS contains everything.'' It does not use the graph or topology in its final form.}

\Q{Population graphs already exist (Parisot et al.). What is new?}
\A Population graphs for disease prediction existed, and in DR they connected images by
visual similarity only. TOPORET adds persistence descriptors of the vessel skeleton to the
edge similarity and to the node features; to our knowledge this had not been done before.
DGTS adds an inductive graph with a feature-space persistent-homology regulariser, and DOTS
adds a screening-oriented evaluation in which safety metrics are reported next to QWK and
checked on external data.

\Q{How do the three systems differ in the evidence they provide?}
\A Table~1 of the Introduction: TOPORET gives transductive evidence (H1, H2), DGTS
inductive and external evidence (EyePACS), DOTS screening-oriented point-estimate
evidence (Messidor-2, H3). They are not identical forms of evidence, and the thesis says so.

\subsection{H1 and TOPORET}

\Q{Your H1 test has four degrees of freedom. Is that serious?}
\A It is a limitation and I treat it as one. The pre-specified test gives $t_4=3.82$,
$p=0.019$. Because five folds are few, I added a bootstrap interval that excludes zero and
an exact sign-flip test ($p=1/32$, the minimum possible: the gain is positive in every fold).
Because folds overlap, I applied the Nadeau--Bengio correction: no single benchmark stays
below $0.05$, but the three benchmarks, which are independent datasets, combined by
Fisher's method give $p=0.015$. So H1 is supported by converging but limited evidence, not
proven in general.

\Q{TOPORET is transductive. Is that not information leakage?}
\A No labels of test images are used; their features, not their grades, influence their
neighbours during training, which is the standard transductive setting of Parisot et al.
It is, however, not equivalent to a strictly inductive, patient-independent evaluation,
and the thesis states it. The inductive evidence for the relational axis comes from DGTS,
whose graph is built from the training fold only (graph $+0.018$ QWK).

\Q{Kaggle DR and Messidor-2 contain both eyes. Could the fellow eye explain the gain?}
\A Partly, possibly: a test image can be connected to the fellow eye, whose grade is
usually the same. Two facts limit this concern: the gain is also observed on APTOS~2019,
which has one image per patient ($+0.04$ QWK), and DGTS, evaluated on a corpus where the
split is patient-level, also shows a graph gain. A patient-level re-evaluation of TOPORET
is the first item of the perspectives; the released code needs only a grouped split.

\Q{Why is your Kaggle DR accuracy not comparable with published systems?}
\A Because it is obtained transductively with image-level folds, which differs from the
patient-level or external test sets of most studies. For that reason Section~4.4.8 does not
rank TOPORET against published systems; what is valid is the paired comparison with the
CNN baseline on the same features and folds.

\Q{Which component of TOPORET matters most?}
\A The graph: with CNN-only node features and topology only in the edges it adds $+0.03$
QWK, against $+0.02$ for topological node features alone; the full model adds $+0.04$.
The contributions are complementary and roughly additive.

\subsection{H2 and topology}

\Q{$\eta^2\approx0.01$: is the effect not negligible?}
\A It is small and I say so. The association is statistically robust -- $F=84.3$ against a
critical value of about $3.72$, and still about $42$ if the effective sample size were
halved -- but grade explains about $1\%$ of the variance of $H_1$ persistence. Topology
alone does not explain severity; it is complementary information, consistent with its
$+0.02$ QWK in the ablation.

\Q{Is $H_1$ persistence a biomarker of neovascularisation?}
\A No. It measures loops of the segmented vessel skeleton. A link with
neovascularisation is plausible but not demonstrated; the thesis uses the term
``severity-associated topological feature'', not ``biomarker''.

\Q{You tested six descriptors. What about multiple comparisons?}
\A The decision level was fixed at $\alpha=0.005$, below the Bonferroni level
$0.05/6\approx0.008$, so the decision on $\Pi_{H_1}$ is valid after correction
(Appendix~C.6).

\Q{Why an ANOVA rather than an ordinal regression?}
\A H2 was formulated as an association between one predefined feature and five grades
treated as groups; the ANOVA makes no assumption on the spacing between grades. Ordinal
or regression models of the grade are natural complementary analyses.

\Q{The topology of TOPORET uses six summary statistics. Is that ``topology''?}
\A It is a coarse summary of the persistence diagrams, and the thesis states that the
topological similarity compares six statistics rather than full diagrams. Richer
representations (persistence images, learnable layers) are listed as perspectives.

\subsection{DGTS}

\Q{What does the persistent-homology branch of DGTS do at inference?}
\A At inference $p_i=0$, so $q_i=\mathrm{ReLU}(W_p p_i)=0$: the topological branch acts only
through the gradients that shaped the graph representation during training. It is a
training-time regulariser. All reported DGTS figures were measured in this inference
setting, so the gain is a gain of the deployed model.

\Q{Is DGTS topology related to the vessels?}
\A No. It is the topology of feature-space neighbourhoods; its loops cannot be read as
vascular loops. The anatomical interpretation of H2 applies only to TOPORET.

\Q{DGTS generalises well but misses the safety threshold. What does that teach?}
\A That accuracy and safety must be measured separately: QWK $0.925$ on EyePACS, yet Sev-NR
about $2.1\%$, slightly above the $2\%$ target. It motivated DOTS.

\subsection{DOTS and H3}

\Q{Is DOTS safe?}
\A Its observed performance satisfies the pre-set screening targets at the point
estimate, internally and on Messidor-2. Safety is not demonstrated: with $309$ severe
images the internal interval is $[0.2\%,\,2.8\%]$, above $2\%$, and the external intervals
cannot be computed from the published counts. About $650$ severe images at the observed
rate would be needed for the upper bound to fall below $2\%$.
\Avoid{``DOTS is safe'' or ``DOTS is clinically validated.''}

\Q{Where does the $2\%$ threshold come from?}
\A It is an evaluation target of this thesis, taken from the DOTS article and informed by
autonomous screening practice (the pivotal trial with pre-defined endpoints) and by
national programmes that monitor the referral of sight-threatening disease. It is not a
regulatory threshold.

\Q{Your non-referral definition counts Severe predicted as Moderate as referred. Is that
lenient?}
\A Yes, and it is stated with the definition. It matches a pathway in which moderate DR is
already referable; a definition requiring a predicted grade of at least~3 would give higher
rates.

\Q{Why is H3b meaningful?}
\A On Messidor-2, every reduced variant fails the Sev-NR target (from $5.8\%$ down to
$2.9\%$), and only the complete pipeline reaches $1.4\%$. The claim is limited to the variants
of this ablation; it does not show that every system lacking a component would fail.

\Q{Your FRE is called a transformer block. Does attention help?}
\A No. It operates on a sequence of length one, so attention weights are identically one:
functionally it is a residual MLP with layer normalisation. Its benefit over a linear
projection comes from the non-linear bottleneck and the residual path, and no
attention-based benefit is claimed.

\Q{CORAL adds only $+0.006$ QWK. Why keep it?}
\A Because the component that improves safety most is not the one that improves QWK most.
CORAL gives the largest accuracy gain per QWK point ($+1.3$~pp), concentrates errors on
adjacent grades ($82.9\%$) and, combined with the other components, is needed to reach the
Sev-NR target externally.

\subsection{Data and protocol}

\Q{KaggleDR-F1 is exactly $60\%$ of APTOS in every grade. Are the sources independent?}
\A Not verified. The proportionality is compatible with KaggleDR-F1 being derived from
APTOS~2019, so the cross-validation results on the fused corpus are interpreted as internal
multi-source evidence, possibly optimistic. The external evaluations (EyePACS, checked for
duplicates, and Messidor-2) are not affected, which is why the conclusions rest on them.

\Q{Are your sample sizes patients or images?}
\A Images. For the fused corpus the split is also patient-level (APTOS has one image per
patient, KaggleDR-F1 no bilateral pairs); for TOPORET on Kaggle DR and Messidor-2 it is not.

\Q{Why a single seed?}
\A The folds were drawn with seed $42$ for reproducibility. Variability over folds is
reported everywhere (Table~D.2), but repetition over several seeds was not run; it is the
second perspective.

\Q{Why EfficientNet and not RETFound or a ViT?}
\A The thesis isolates the contribution of relations, topology and ordinality on a fixed,
widely used backbone, so the gains are not confounded by a change of backbone. Retinal
foundation models such as RETFound are listed as a methodological perspective.

\Q{No demographic information: what about fairness?}
\A None of the public datasets provides per-image demographic metadata, so a subgroup
analysis could not be done; it is stated as a limitation and is a requirement of the
planned prospective study.

\subsection{Chapter 3 and the non-medical detour}

\Q{Why evaluate a graph template on microservices and network intrusion?}
\A Those domains have unambiguous labels and graphs with a clear structure, so the template
could be tested before being applied to patients. The ten-run re-evaluation shows an
honest result: the invariants bring interpretability and robustness to drift, not a
significant accuracy gain. The transfer to retinal grading is a design analogy, not a
validated result.

\Q{What does SAMF-GB show?}
\A That structural descriptors alone already carry enough information to rival a deep
network on this task ($+0.042$ QWK over EfficientNet-B3, with a CPU classifier about $13$
times cheaper). It is not a like-for-like architecture benchmark.

\subsection{Chapter 6 and deployment}

\Q{Is your system compliant with the MDR or the EU AI Act?}
\A No claim of compliance is made. Chapter~6 is a regulatory-oriented engineering mapping,
not a conformity assessment, clinical evaluation, CE marking or FDA clearance. Table~6.6
states the status of each component: the FHIR adapter is designed but not implemented, and
no deployment or certification was performed.

\Q{What would the next study look like?}
\A A pre-registered, prospective, multi-site reader study powered on severe cases: about
$1{,}240$ severe images for $80\%$ power, i.e.\ about $25{,}000$ screened images at a $5\%$
prevalence of severe disease, with patient-level splits and a deferral mechanism based on
the calibrated probabilities of DOTS.

\subsection{Trap questions}

\Q{So what did you really prove?}
\A The experiments show directly the measured quantities under the stated protocols. They
support, without establishing conclusively, that relational information improves grading,
that topology carries severity-related information, and that the ordinal pipeline
satisfies screening targets. They do not evaluate clinical safety, patient-level TOPORET,
multi-seed variability, fairness, longitudinal prediction or regulatory compliance
(Table~GC.2).

\Q{If you had six more months, what would you do first?}
\A Re-run TOPORET with patient-level folds and a per-fold graph, repeat all experiments over
several seeds, and add a second external dataset with enough severe cases. These three
steps address the main statistical limits without new infrastructure.

\Q{The title says ``disease severity''. Do you predict progression?}
\A No. The task is cross-sectional: severity on the ICDR scale from a single image. A
temporal population graph linking successive visits is a perspective.

\section{Wording to use and to avoid}
\begin{center}\small
\begin{tabular}{@{}P{7.6cm}P{8.2cm}@{}}
\toprule
\color{limred}Avoid & \color{okgreen}Say instead \\
\midrule
``We proved that the graph improves grading'' & ``The experiments provide evidence that relational information improves grading under the evaluated conditions'' \\
``Topology explains severity'' & ``A small statistical association between a topological descriptor and severity'' \\
``DOTS is safe'' & ``Its observed performance satisfies the pre-set screening targets at the point estimate'' \\
``Clinically validated'' & ``Evaluated retrospectively on external public data'' \\
``Regulatory-compliant'' & ``Designed with regulatory considerations in mind'' \\
``Patients'' (for dataset counts) & ``Retinal images'' \\
``Retinal topology'' (for DGTS) & ``Feature-space topological regularisation'' \\
``Superior to the state of the art'' & ``Higher than the baseline on the same folds and features'' \\
\bottomrule
\end{tabular}
\end{center}

\section{Suggested presentation plan (about 40 minutes)}
\begin{center}\small
\begin{tabular}{@{}rP{4.8cm}P{9.2cm}r@{}}
\toprule
\# & Part & Content & min \\
\midrule
1 & Context & DR burden, screening gap, why accuracy is not enough & 3 \\
2 & Problem & L1--L3; the three hypotheses; objective & 3 \\
3 & Methodology & Pre-specified tests; framework of three systems (Table~1) & 3 \\
4 & Foundations & Chapter~3 in one slide: template outside medicine, structural evidence & 3 \\
5 & TOPORET & Dual-similarity graph; H1 with its robustness and correction; ablation & 7 \\
6 & H2 & ANOVA, effect size, what it does not mean & 3 \\
7 & DGTS & Inductive graph, regulariser at training only, EyePACS, Sev-NR $2.1\%$ & 5 \\
8 & DOTS & Pipeline, ablation, H3a with interval, H3b on Messidor-2 & 6 \\
9 & Engineering & Six stages, status table, what is not done & 2 \\
10 & Evidence and limits & Figure~GC.1 and Table~GC.2 & 3 \\
11 & Perspectives & Patient-level, multi-seed, external data, reader and prospective study & 2 \\
\bottomrule
\end{tabular}
\end{center}
Open and close with the same sentence: \emph{relational, topological and ordinal
information can complement deep visual representations for DR grading -- here is the
evidence, and here is exactly where it stops.}

\end{document}
